
\section{\BookRecordIndex{T0244}冒泡现象}\label{sec:4-2}
\zhindex{maopaoxianxiang}{冒泡现象}

设 $(M,\omega)$ 是紧\BookRecordIndex{T0358}辛流形，$J\in\mathcal J_\tau(M,\omega)$ 是 $M$ 上被
$\omega$ \BookRecordIndex{T0372}驯服的\BookRecordIndex{T0088}殆复结构，并设
$(\Sigma,j_\Sigma,d\operatorname{vol}_\Sigma)$ 是闭黎曼曲面。我们在
第~\ref{sec:4-1} 节已经看到，一阶导数一致有界的任意 \BookRecordIndex{T0003}$J$-全纯曲线列
$u^\nu:\Sigma\to M$ 都有一个 $C^\infty$ 收敛子列（定理~\ref{thm:4-1-1}）。
因此，只有当映射列 $du^\nu$ 无界时，\BookRecordIndex{T0192}紧性才可能失效：
\[
  \sup_\nu\lVert du^\nu\rVert_{L^\infty}=\infty.
\]
若能量还一致有界，即
\[
  \sup_\nu E(u^\nu)<\infty,
\]
那么\BookRecordIndex{T0148}共形重标度论证表明，会有一个全纯球面冒出。下面说明这一过程。

令 $z^\nu\in\Sigma$ 是实值函数 $|du^\nu|$ 取得最大值的一点：
\[
  |du^\nu(z^\nu)|
  =\lVert du^\nu\rVert_{L^\infty(\Sigma)}=:c^\nu.
\]
必要时取子列，可设 $z^\nu$ 收敛到一点 $z_0\in\Sigma$，而 $c^\nu$ 趋于无穷。
现在取定义在零点的一个开邻域 $\Omega\subset\mathbb C$ 上的全纯坐标卡
$\phi:\Omega\to\Sigma$，使得 $\phi(0)=z_0$。于是拉回体积形式为
\[
  \phi^*d\operatorname{vol}_\Sigma=\lambda^2\,ds\wedge dt,
\]
其中 $\lambda:\Omega\to(0,\infty)$。不妨假设对每个 $\nu$，点 $z^\nu$ 都落在
坐标卡的像 $\phi(\Omega)$ 中\sourcefootnote{原书此处写作 $z^\nu\in\Omega$；但此前
$z^\nu\in\Sigma$，而坐标域 $\Omega\subset\mathbb C$，故据坐标卡
$\phi:\Omega\to\Sigma$ 的源靶改作 $z^\nu\in\phi(\Omega)$。}，并且对所有
$z\in\Omega$，
\[
  \lambda(0)=1,
  \qquad
  \frac12\le\lambda(z)\le2.
\]
考虑映射列与点列
\[
  u^\nu_{\mathrm{loc}}:=u^\nu\circ\phi:\Omega\to M,
  \qquad
  z^\nu_{\mathrm{loc}}:=\phi^{-1}(z^\nu).
\]
以下略去下标 ``$\mathrm{loc}$''。于是
\[
  c^\nu
  =\frac{|du^\nu(z^\nu)|}{\lambda(z^\nu)}
  =\sup_\Omega\frac{|du^\nu|}{\lambda},
  \qquad
  \lim_{\nu\to\infty}z^\nu=0.
\]

现在选取 $\varepsilon>0$，使得对每个 $\nu$ 都有
$B_\varepsilon(z^\nu)\subset\Omega$，并考虑重参数化映射列
$v^\nu:B_{\varepsilon c^\nu}\to M$，其定义为
\[
  v^\nu(z):=u^\nu\bigl(z^\nu+z/c^\nu\bigr).
\]
这个映射列满足
\[
  |dv^\nu(0)|\ge\frac12,
  \qquad
  \lVert dv^\nu\rVert_{L^\infty(B_{\varepsilon c^\nu})}\le2,
\]
以及
\[
  E(v^\nu;B_{\varepsilon c^\nu})
  =E\bigl(u^\nu;B_\varepsilon(z^\nu)\bigr)
  \le E(u^\nu),
\]
其中 $E(u;B)$ 表示 $u$ 限制在 $B$ 上的能量。最后一个恒等式来自能量的共形不变性。
由定理~\ref{thm:4-1-1}，存在一个子列（仍记作 $v^\nu$），它在紧子集上一致且连同
所有阶导数一起收敛。极限映射 $v:\mathbb C\to M$ 仍是 \BookRecordIndex{T0003}$J$-全纯曲线，并满足
\[
  |dv(0)|\ge\frac12,
  \qquad
  0<E(v)=\int_{\mathbb C}v^*\omega
  \le\sup_\nu E(u^\nu).
\]
能量的共形不变性说明映射
\[
  \mathbb C\setminus\{0\}\longrightarrow M,
  \qquad z\longmapsto v(1/z),
\]
具有有限能量；所以由\BookRecordIndex{T0214}可去奇点定理~\ref{thm:4-1-2}，它可以光滑延拓过零点。
因此，$v$ 延拓为从黎曼球面 $S^2\cong\mathbb C\cup\{\infty\}$ 到 $M$ 的非常值
$J$-全纯映射。以这种方式构造出的 $J$-全纯球面 $v$ 称为一个\emph{\BookRecordIndex{T0275}气泡}。
\zhindex{qipao}{气泡}
这个\BookRecordIndex{T0275}气泡的能量 $E(v)$ 为正，而且当 $\nu\to\infty$ 时，$u^\nu$ 在点 $z_0$ 的任意小
邻域中的能量至少为 $E(v)$；也就是说，对每个 $\varepsilon>0$，都有
\begin{equation}\label{eq:4-2-1}
  \liminf_{\nu\to\infty}E\bigl(u^\nu;B_\varepsilon(z_0)\bigr)\ge E(v).
\end{equation}
事实上，对每个 $R>0$ 和每个 $\varepsilon>0$，下面推导的最后一行应取极限下限
\sourcefootnote{原书最后一行写作 $\le\lim_{\nu\to\infty}$，但并未建立右侧
能量列的极限存在；这里的目标式~\eqref{eq:4-2-1} 使用 $\liminf$，且由逐项包含关系
唯一得到的也是 $\liminf$，故据此补正。}：
\[
  \begin{aligned}
  E(v;B_R)
  &=\lim_{\nu\to\infty}E(v^\nu;B_R)\\
  &=\lim_{\nu\to\infty}
    E\bigl(u^\nu;B_{R/c^\nu}(z^\nu)\bigr)\\
  &\le\liminf_{\nu\to\infty}
    E\bigl(u^\nu;B_\varepsilon(z_0)\bigr).
  \end{aligned}
\]
令 $R\to\infty$，即得不等式~\eqref{eq:4-2-1}。

\paragraph{\BookRecordIndex{T0040}Möbius 变换群。}
上述重标度论证的关键是能量具有共形不变性。最早可以显式观察\BookRecordIndex{T0244}冒泡现象的例子之一，是
黎曼球面的共形自同构群
$G:=\operatorname{PSL}(2,\mathbb C)$，它也称为 \BookRecordIndex{T0040}Möbius 变换群。其元素是\BookRecordIndex{T0126}分式线性变换
\begin{equation}\label{eq:4-2-2}
  \phi(z)=\frac{az+b}{cz+d},
  \qquad ad-bc=1,
\end{equation}
其中 $a,b,c,d\in\mathbb C$。这些恰好是保持复结构的二维球面微分同胚；因而，如第 3 章
所述，它们组成从 $S^2$ 到 $M=S^2$ 的一次全纯映射之\BookRecordIndex{T0246}模空间。

群 $G$ 不是紧群；它的一个极大紧子群是 $S^2$ 的等距变换群
$\operatorname{SO}(3)=\operatorname{SU}(2)/\{\pm\mathbf1\}$。$G$ 的非\BookRecordIndex{T0192}紧性与能量的
共形不变性有关，并可用\BookRecordIndex{T0243}冒泡来解释。

\paragraph{例与习题。}

\begin{exercise}\label{ex:4-2-1}
利用从北极出发的球极投影
\[
  S^2\longrightarrow\mathbb C\cup\{\infty\},
  \qquad
  (x_1,x_2,x_3)\longmapsto\frac{x_1+ix_2}{1-x_3},
\]
把单位球面 $S^2\subset\mathbb R^3$ 与黎曼球面 $\mathbb C\cup\{\infty\}$ 等同。
证明 Fubini--Study 形式
\[
  (1+x^2+y^2)^{-2}\,dx\wedge dy,
  \qquad z=x+iy,
\]
的拉回，等于 $S^2$ 上标准面积形式
\[
  x_1\,dx_2\wedge dx_3+x_2\,dx_3\wedge dx_1+x_3\,dx_1\wedge dx_2
\]
的负四分之一。（这里的负号是必要的，因为 $S^2$ 作为 $\mathbb R^3$ 中一个球的边界
所带的定向，与它作为 $\mathbb C\cup\{\infty\}$ 所带的定向不同。）
\end{exercise}

\begin{exercise}\label{ex:4-2-2}
设 $\phi$ 是形如式~\eqref{eq:4-2-2} 的\BookRecordIndex{T0126}分式线性变换。证明 $\phi$ 保持上半平面
\[
  \mathbb H:=\{z\in\mathbb C\mid\operatorname{Im}z\ge0\}
\]
当且仅当 $a,b,c,d\in\mathbb R$。证明\BookRecordIndex{T0126}分式线性变换
\[
  \phi(z):=\frac{i-z}{i+z}
\]
把上半平面等同于 $B\setminus\{-1\}$，其中
\[
  B:=\{z\in\mathbb C\mid |z|\le1\}.
\]
其逆映射为
\[
  B\setminus\{-1\}\longrightarrow\mathbb H,
  \qquad z\longmapsto i(1-z)(1+z)^{-1}.
\]
\end{exercise}

\begin{exercise}\label{ex:4-2-3}
形如式~\eqref{eq:4-2-2} 的 \BookRecordIndex{T0040}Möbius 变换 $\phi$ 的导数为
\[
  \phi'(z)=\frac1{(cz+d)^2},
\]
而关于 $S^2=\mathbb C\cup\{\infty\}$ 上 Fubini--Study 度量的导数范数为
\[
  \frac1{\sqrt2}|d\phi(z)|_{\mathrm{FS}}
  =|\phi'(z)|\frac{1+|z|^2}{1+|\phi(z)|^2}
  =\frac{1+|z|^2}{|az+b|^2+|cz+d|^2}.
\]
（这里定义域和靶上都使用 Fubini--Study 度量；另见习题~\ref{ex:2-2-3}。）不用
定理~\ref{thm:4-1-1}，直接证明：每个满足
\[
  \sup_\nu\sup_z|d\phi^\nu(z)|_{\mathrm{FS}}<\infty
\]
的映射列 $\phi^\nu\in G$，都有一个在整个 $S^2$ 上一致且连同所有阶导数一起收敛的
子列。
\end{exercise}

\begin{exercise}\label{ex:4-2-4}
使用上面概述的\BookRecordIndex{T0243}冒泡论证，给出引理 D.1.2 的另一个证明：若 $\phi^\nu\in G$ 是一列
没有一致收敛子列的 \BookRecordIndex{T0040}Möbius 变换，则存在点 $x,y\in S^2$ 和一个子列
$\phi^{\nu_i}$，使它在 $S^2\setminus\{x\}$ 的紧子集上一致收敛到 $y$。
提示：若遇到困难，可参看定理~\ref{thm:4-6-1}，并利用如下事实：对每个
$\phi\in G$，
\[
  E(\phi)=\pi=\hbar(S^2,\omega_{\mathrm{FS}},J_{\mathrm{FS}}).
\]
\end{exercise}

\begin{exercise}\label{ex:4-2-5}
双全纯映射 $\phi:B\to B$ 是形如
\[
  \phi(z)=e^{i\theta}\frac{z-z_0}{1-\bar z_0z}
\]
的\BookRecordIndex{T0126}分式线性变换，其中 $\theta\in\mathbb R$ 且 $z_0\in\operatorname{int}(B)$。证明
\[
  |\phi'(z)|=\frac{1-|\phi(z)|^2}{1-|z|^2}.
\]
由此\BookRecordIndex{T0326}推出：若在定义域和靶上都使用双曲度量
$(dx^2+dy^2)/(1-x^2-y^2)^2$，则 $\phi$ 是等距映射。再考虑映射列
\[
  \phi^\nu(z):=\frac{z-x^\nu}{1-x^\nu z},
\]
其中 $x^\nu\in(0,1)$ 且 $x^\nu\to1$。证明 $(\phi^\nu)'(1)$ 趋于无穷，而对每个
$z\ne1$，$(\phi^\nu)'(z)$ 都趋于零。用具有拉格朗日边界条件的 $J$-全纯圆盘在边界上
\BookRecordIndex{T0243}冒泡来解释这种极限行为。
\end{exercise}

\begin{exercise}\label{ex:4-2-6}
考虑 $M=\Sigma=S^2=\mathbb C\cup\{\infty\}$ 的情形。全纯曲线是一个有理函数
\[
  u(z)=\frac{p(z)}{q(z)}.
\]
若 $p$ 与 $q$ 互素，则 $u$ 的次数是分子与分母次数的较大者。与
习题~\ref{ex:4-2-3} 一样，它关于 Fubini--Study 度量的导数范数为
\[
  \frac1{\sqrt2}|du(z)|_{\mathrm{FS}}
  =|u'(z)|\frac{1+|z|^2}{1+|u(z)|^2}.
\]
现在考虑一列 $k+1$ 次有理映射
\[
  u^\nu(z):=\frac{zp(z)}{(z-a^\nu)q(z)}.
\]
假设 $p$ 与 $q$ 互素且 $p(0)=q(0)=1$，并假设 $a^\nu\to0$。证明在 $z=0$ 处冒出
一个一次全纯球面，而且 $u^\nu$ 在 $S^2\setminus\{0\}$ 的紧子集上一致且连同所有
阶导数一起收敛到 $u=p/q$。
\end{exercise}

\begin{exercise}\label{ex:4-2-7}
考虑全纯映射列 $u^\nu:\mathbb{CP}^1\to\mathbb{CP}^2$，其定义为
\[
  u^\nu([x:y]):=[x^2:y^2:\nu xy].
\]
其像是二次曲线族 $XY=Z^2/\nu^2$；当 $\nu\to\infty$ 时，它退化为两条直线
$X=0$ 与 $Y=0$。适当重参数化给定映射列 $u^\nu$，把这两条极限曲线分别实现为所形成
的\BookRecordIndex{T0275}气泡。
\end{exercise}

\begin{exercise}\label{ex:4-2-8}
对映射列 $u^\nu:S^2\to S^2\times S^2$ 考察与习题~\ref{ex:4-2-7} 相同的问题，其中
\[
  u^\nu(z):=(z,\phi^\nu(z)),
\]
而 $\phi^\nu$ 如习题~\ref{ex:4-2-4} 所述。
\end{exercise}

\begin{example}\label{ex:4-2-9}
考虑 $\Sigma=T^2=\mathbb C/\Lambda$，其中 $\Lambda=\mathbb Z+i\mathbb Z$。
Weierstrass $p$-函数是如下定义的亚纯函数 $p:T^2\to S^2$：
\[
  p(z):=\sum_{\lambda\in\Lambda}\frac1{(z-\lambda)^2}.
\]
这个映射是满射，并且
$p(z_1)=p(z_2)$ 当且仅当 $z_1\pm z_2\in\Lambda$，所以它的次数为 $2$。又有
$c_1(p^*TS^2)=4$，故从 $T^2$ 到 $S^2$ 的二次亚纯映射之\BookRecordIndex{T0246}模空间 $\mathcal M_2$
具有实维数 $8$。每个 $u\in\mathcal M_2$ 都可以写成
\[
  u(z)=\phi\bigl(p(z+z_0)\bigr),
\]
其中 $\phi\in\operatorname{PSL}(2,\mathbb C)$ 且 $z_0\in\mathbb C$。考虑映射列
\[
  u^\nu(z):=\frac{p(z)-p(z_0)}{p(z)-a^\nu p(z_0)},
\]
其中 $p(z_0)$ 是 $p$ 的一个非零\BookRecordIndex{T0409}正则值。若 $a^\nu\to1$，则 $u^\nu$ 的导数在
$\pm z_0$ 附近趋于无穷：
\[
  (u^\nu)'(\pm z_0)
  =\frac{p'(\pm z_0)}{p(z_0)(1-a^\nu)}.
\]
因此，在 $\pm z_0$ 处分别冒出一个一次全纯球面，而 $u^\nu$ 在
$T^2\setminus\{-z_0,z_0\}$ 的每个紧子集上以 $C^\infty$ 拓扑收敛到常值 $1$。
\end{example}

\begin{remark}\label{rem:4-2-10}
当定义域 $\Sigma$ 具有正亏格时，若允许 $\Sigma$ 上的复结构 $j$ 变化，还可能出现
进一步的退化。可以显式观察这一现象的第一个例子是 $\mathbb{CP}^2$ 中的三次曲线。
一般的三次曲线是嵌入环面；它们例如可以退化为有一个自交点或一个奇点的\BookRecordIndex{T0386}有理曲线，或
三重覆盖的直线。一个可约极限的例子是三条直线之并。关于这种极限行为的细致研究，见
Lizan \,\srcite{252}。
\end{remark}

\paragraph{具有最小能量的球面。}
若同调类 $A\in H_2(M)$ 位于胡雷维茨同态
$\pi_2(M)\to H_2(M)$ 的像中，就称它为\emph{\BookRecordIndex{T0285}球面类}。\BookRecordIndex{T0285}球面类的集合记作
$H_2^S(M)$。固定球面同调类 $A\in H_2(M;\mathbb Z)$，使得
\[
  \langle[\omega],A\rangle=\hbar,
\]
其中 $\hbar=\hbar(M,\omega,J)$ 是命题~\ref{prop:4-1-4} 的常数。以
$\mathcal M(A;J)$ 表示代表类 $A$ 的 $J$-全纯球面 $u:S^2\to M$ 之\BookRecordIndex{T0246}模空间。由\BookRecordIndex{T0214}可去
奇点定理，我们可以把全纯球面看成一条有限能量全纯曲线 $u:\mathbb C\to M$。于是映射
\[
  \mathbb C\setminus\{0\}\longrightarrow M,
  \qquad z\longmapsto u(1/z),
\]
可以光滑延拓过零点，所以 $u$ 延拓为从黎曼球面
$\mathbb C\cup\{\infty\}$ 到 $M$ 的 $J$-全纯映射。\BookRecordIndex{T0040}Möbius 变换群
$G=\operatorname{PSL}(2,\mathbb C)$ 通过重参数化作用在 $\mathcal M(A;J)$ 上。
若 \BookRecordIndex{T0003}$J$-全纯曲线列 $u^\nu:\mathbb C\to M$ 中的 $u^\nu(z)$ 与 $u^\nu(1/z)$ 都在
$\mathbb C$ 的紧子集上一致且连同所有阶导数一起收敛，就称 $u^\nu$ 在 $S^2$ 上以
$C^\infty$ 拓扑收敛。

\begin{theorem}\label{thm:4-2-11}
设 $(M,\omega)$ 是紧\BookRecordIndex{T0358}辛流形，并固定 $J\in\mathcal J_\tau(M,\omega)$。若
$A\in H_2(M;\mathbb Z)$ 是满足
$\langle[\omega],A\rangle=\hbar$ 的球面同调类，则\BookRecordIndex{T0246}模空间
$\mathcal M(A;J)/G$ 是紧的。换言之，对代表类 $A$ 的每一列 $J$-全纯球面
$u^\nu:S^2\to M$，都存在一列 $\phi^\nu\in G$，使 $u^\nu\circ\phi^\nu$ 有一个
$C^\infty$ 收敛子列。
\end{theorem}

\begin{proof}
先复合一列 $\operatorname{SO}(3)$ 中的 \BookRecordIndex{T0040}Möbius 变换，不妨假设对每个 $\nu$ 都有
\[
  |du^\nu(0)|_{\mathrm{FS}}
  =\sup_{z\in\mathbb C}|du^\nu(z)|_{\mathrm{FS}}.
\]
由于
\[
  |du^\nu(z)|_{\mathrm{FS}}
  =|du^\nu(z)|(1+|z|^2),
\]
这蕴含
\[
  |du^\nu(0)|=\sup_{z\in\mathbb C}|du^\nu(z)|.
\]
再作一个以零点为中心的\BookRecordIndex{T0148}共形重标度，不妨进一步假设
\[
  |du^\nu(0)|=1.
\]
由定理~\ref{thm:4-1-1}，$u^\nu$ 有一个子列（仍记作 $u^\nu$），它在
$\mathbb C$ 的紧子集上一致且连同所有阶导数一起收敛到一条非常值 \BookRecordIndex{T0003}$J$-全纯曲线
$u:\mathbb C\to M$。现在考虑映射列 $z\mapsto u^\nu(1/z)$。它在
$\mathbb C\setminus\{0\}$ 的紧子集上一致且连同所有阶导数一起收敛到 $u(1/z)$。
我们断言，$u^\nu(1/z)$ 在零点附近也一致收敛。否则就会发生\BookRecordIndex{T0243}冒泡，而由
式~\eqref{eq:4-2-1}，对每个 $R>0$ 都有
\[
  \liminf_{\nu\to\infty}E(u^\nu;\mathbb C\setminus B_R)\ge\hbar.
\]
这将\BookRecordIndex{T0326}推出 $E(u;B_R)=0$，与 $u$ 非常值矛盾。因此，$u^\nu$ 在整个 $S^2$ 上以
$C^\infty$ 拓扑收敛到 $u$。
\end{proof}

本节的结果已经足以得到辛拓扑中 \BookRecordIndex{T0003}$J$-全纯曲线的若干卓越应用。这些应用由
\personlater{geluomofu}{格罗莫夫}{Gromov}在 \,\srcite{160} 中给出，并不需要
第 5 章更精细的\BookRecordIndex{T0192}紧性结果（见第~\ref{ch:applications-symplectic-topology} 章）；它们使我们初步领略到 \BookRecordIndex{T0003}$J$-全纯曲线理论在
辛拓扑中的威力。
